\begin{abstract}
\noindent Manual font design is an intricate process that transforms a stylistic visual concept into a coherent glyph set. This challenge persists in automated Few-shot Font Generation (FFG), where models often struggle to preserve both the structural integrity and stylistic fidelity from limited references. While autoregressive (AR) models have demonstrated impressive generative capabilities, their application to FFG is constrained by conventional patch-level tokenization, which neglects global dependencies crucial for coherent font synthesis. Moreover, existing FFG methods remain within the image-to-image paradigm, relying solely on visual references and overlooking the role of language in conveying stylistic intent during font design. To address these limitations, we propose GAR-Font, a novel AR framework for multimodal few-shot font generation. GAR-Font introduces a global-aware tokenizer that effectively captures both local structures and global stylistic patterns, a multimodal style encoder offering flexible style control through a lightweight language-style adapter without requiring intensive multimodal pretraining, and a post-refinement pipeline that further enhances structural fidelity and style coherence. Extensive experiments show that GAR-Font outperforms existing FFG methods, excelling in maintaining global style faithfulness and achieving higher-quality results with textual stylistic guidance.


% \noindent\chn{Manual font design is an intricate process that transforms a stylistic vision into a consistent glyph set. This challenge persists in automated Few-shot Font Generation (FFG), where models often struggle to preserve both the structural integrity and stylistic consistency from limited references. While autoregressive (AR) models have demonstrated impressive generative capabilities, their application to FFG is constrained by conventional patch-level tokenization approaches, which neglect global dependencies crucial for coherent font synthesis. Moreover, existing FFG methods remain within the image-to-image paradigm, relying solely on visual references and overlooking the role of language in conveying stylistic intent during font design. To address these limitations, we propose GAR-Font, a novel AR framework for multimodal few-shot font generation. GAR-Font introduces a global-aware tokenizer that effectively captures both local structures and global stylistic patterns, and a lightweight language-style adapter that provides multimodal stylistic guidance to improve generation quality without intensive multimodal pretraining. Through a combination of supervised fine-tuning and reinforcement learning, GAR-Font is further refined for superior performance. Extensive experiments show that GAR-Font outperforms existing few-shot font generation methods, excelling in maintaining global style consistency across generated glyphs and achieving higher-quality results when textual cues complement few-shot visual inputs.}%Lian
\end{abstract}